../cictikz/src/cictikz/data/tex/cictikz_preamble.tex

% Ring oscillator frequency and its slope against temperature.
%
% Frequency falls by a factor of 2.6 from -40 to 150 degrees, because
% mobility falls with temperature and the inverters get slower. A ring
% oscillator is therefore a usable temperature sensor and an unusable
% clock, and the two statements are the same statement.
%
% The lower panel gives the number to design with: roughly -12 MHz per
% degree in the cold, -3 MHz per degree when hot. The sensitivity is
% itself temperature dependent, which is what makes compensating a ring
% oscillator harder than it first looks.
%
% Generated by py/tikzplot.py. Edit the script in ex/ that
% produced it and re-run that, not this file.

\begin{tikzpicture}[thick]

\begin{groupplot}[
  group style={group size=1 by 2, horizontal sep=1.7cm, vertical sep=1.5cm},
  scale only axis,
  grid=both,
  grid style={gray!22, very thin},
  tick align=outside,
  tick label style={font=\small},
  label style={font=\small},
  title style={font=\small, yshift=-1mm},
  legend cell align=left,
  legend style={font=\small, draw=gray!50, fill=white, fill opacity=0.85, text opacity=1}
]
\nextgroupplot[
  width=9.5cm,
  height=3.8cm,
  xlabel={Temperature [$^\circ$C]},
  ylabel={Frequency [MHz]}
]
\addplot[black, very thick, mark=none] coordinates {(-40,2143.6) (-30,2021.1) (-20,1906.2) (-10,1799.3) (0,1699.8) (10,1607.6) (20,1522.2) (30,1443) (40,1369.7) (50,1301.7) (60,1238.7) (70,1180.3) (80,1126.1) (90,1075.7) (100,1028.8) (110,984.9) (120,944.49) (130,906.49) (140,871.26) (150,838.37)};
\nextgroupplot[
  width=9.5cm,
  height=3.8cm,
  xlabel={Temperature [$^\circ$C]},
  ylabel={$df/dT$ [MHz/$^\circ$C]}
]
\addplot[red, very thick, mark=none] coordinates {(-40,-12.246) (-30,-11.868) (-20,-11.093) (-10,-10.319) (0,-9.584) (10,-8.8825) (20,-8.2295) (30,-7.6255) (40,-7.063) (50,-6.5465) (60,-6.071) (70,-5.6325) (80,-5.231) (90,-4.867) (100,-4.5396) (110,-4.2133) (120,-3.9202) (130,-3.6614) (140,-3.406) (150,-3.2884)};
\end{groupplot}

\end{tikzpicture}
\end{document}
